% !PRINTED 324
\weberpage{324}

%<<CONT
wenn $n$ eine ungerade Primzahl ist (nach dem Reziprozitätsgesetz und seinen Ergänzungssätzen [Bd. I, § 145, 9.]), und diese Formeln lassen sich zusammenfassen in
\weq{10}{(\pm p, n) = \left(\dfrac{n}{p}\right).}
Diese Formel gilt aber wegen (7) für \gesperrt{jedes beliebige $n$, das durch $p$ nicht teilbar ist.}

In gleicher Weise findet man für ein ungerades $n$:
\weq{11}{\begin{aligned} (-4,n) &= (-1)^{\frac{n-1}{2}},\\ (8,n) &= (-1)^{\frac{n^2-1}{8}},\\ (-8,n) &= (-1)^{\frac{n-1}{2}+\frac{n^2-1}{8}}. \end{aligned}}
Weiter ist allgemein, wenn $D$, $D'$ zwei Diskriminanten sind:
\weq{12}{(D,n)(D',n) = (DD',n).}
Diese Formel gilt zunächst, wenn $DD'$ und $n$ nicht relativ prim sind, weil dann beide Seiten $=0$ sind. Sie gilt ferner für $n=-1$ [nach (3)], für $n=2$ [nach (5) und der Kongruenz $\dfrac{D^2-1}{8}+\dfrac{D'^2-1}{8} \equiv \dfrac{D^2D'^2-1}{8}$ (mod 2)] und für eine ungerade Primzahl $n$ [nach (6)]. Also gilt sie wegen (7) allgemein.

Wir können nun das Reziprozitätsgesetz mit seinen Ergänzungssätzen für Diskriminanten in folgender Weise zusammenfassen:
\weq{13}{(D,D') = \pm(D',D),}
\gesperrt{worin das obere Zeichen gilt, wenn von den beiden Diskriminanten $D$, $D'$ wenigstens eine positiv ist, das untere, wenn beide negativ sind.}

Um diese Formel aus der gewöhnlichen Form des Reziprozitätsgesetzes abzuleiten, nehmen wir zunächst an, es gelten die beiden Formeln:
\[
(D,D') = \pm(D',D),
\]
\[
(D_1,D') = \pm(D',D_1).
\]
Dann folgt aus (8) und (12)
\[
(DD_1,D') = \pm(D',DD_1),
\]
worin wieder das untere Zeichen nur dann gilt, wenn $D'$ und $DD_1$ beide negativ sind. Durch nochmalige Anwendung ergibt sich
\[
(DD_1,D'D_1') = \pm(D'D_1',DD_1)
\]
%CONT>>
